% Options for packages loaded elsewhere
\PassOptionsToPackage{unicode}{hyperref}
\PassOptionsToPackage{hyphens}{url}
\documentclass[
  ukrainian,
]{article}
\usepackage{xcolor}
\usepackage[margin=2.5cm]{geometry}
\usepackage{amsmath,amssymb}
\setcounter{secnumdepth}{-\maxdimen} % remove section numbering
\usepackage{iftex}
\ifPDFTeX
  \usepackage[T1]{fontenc}
  \usepackage[utf8]{inputenc}
  \usepackage{textcomp} % provide euro and other symbols
\else % if luatex or xetex
  \usepackage{unicode-math} % this also loads fontspec
  \defaultfontfeatures{Scale=MatchLowercase}
  \defaultfontfeatures[\rmfamily]{Ligatures=TeX,Scale=1}
\fi
\usepackage{lmodern}
\ifPDFTeX\else
\setmainfont{DejaVu Serif}
\setsansfont{DejaVu Sans}
\setmonofont{DejaVu Sans Mono}
\setmathfont{STIX Two Math}
\fi
\ifPDFTeX\else
  % xetex/luatex font selection
\fi
% Use upquote if available, for straight quotes in verbatim environments
\IfFileExists{upquote.sty}{\usepackage{upquote}}{}
\IfFileExists{microtype.sty}{% use microtype if available
  \usepackage[]{microtype}
  \UseMicrotypeSet[protrusion]{basicmath} % disable protrusion for tt fonts
}{}
\makeatletter
\@ifundefined{KOMAClassName}{% if non-KOMA class
  \IfFileExists{parskip.sty}{%
    \usepackage{parskip}
  }{% else
    \setlength{\parindent}{0pt}
    \setlength{\parskip}{6pt plus 2pt minus 1pt}}
}{% if KOMA class
  \KOMAoptions{parskip=half}}
\makeatother
\usepackage{graphicx}
\graphicspath{{/Users/kalinicenkovitalijmikolajovic/.codex/worktrees/5c15/Script_for_translate_book/Output_Final/}}
\setkeys{Gin}{width=0.82\linewidth,height=0.82\textheight,keepaspectratio}
\makeatletter
\newsavebox\pandoc@box
\newcommand*\pandocbounded[1]{% scales image to fit in text height/width
  \sbox\pandoc@box{#1}%
  \Gscale@div\@tempa{\textheight}{\dimexpr\ht\pandoc@box+\dp\pandoc@box\relax}%
  \Gscale@div\@tempb{\linewidth}{\wd\pandoc@box}%
  \ifdim\@tempb\p@<\@tempa\p@\let\@tempa\@tempb\fi% select the smaller of both
  \ifdim\@tempa\p@<\p@\scalebox{\@tempa}{\usebox\pandoc@box}%
  \else\usebox{\pandoc@box}%
  \fi%
}
% Set default figure placement to htbp
\def\fps@figure{htbp}
\makeatother
\ifLuaTeX
\usepackage[bidi=basic,shorthands=off,provide=*]{babel}
\else
\usepackage[bidi=default,shorthands=off,provide=*]{babel}
\fi
\ifLuaTeX
  \usepackage{selnolig} % disable illegal ligatures
\fi
\setlength{\emergencystretch}{3em} % prevent overfull lines
\providecommand{\tightlist}{%
  \setlength{\itemsep}{0pt}\setlength{\parskip}{0pt}}
\usepackage{bookmark}
\IfFileExists{xurl.sty}{\usepackage{xurl}}{} % add URL line breaks if available
\urlstyle{same}
\hypersetup{
  pdflang={uk},
  hidelinks,
  pdfcreator={LaTeX via pandoc}}

\author{}
\date{}

\begin{document}

\pandocbounded{\includegraphics[keepaspectratio]{images/_page_273_Picture_10.jpeg}}

36 1 Радиальный поток тепла от точечного источника; a ---тепло
выделяется в начале координат.

\begin{verbatim}
u(r) = \frac{-q}{2\pi} \ln r = \frac{q}{2\pi} \ln \frac{1}{r}.
\end{verbatim}

(Решение изображено на рис. 362)В электростатике разность потенциа тов
\texttt{u(B)} - \texttt{u(A)} равна работе, которую необходимо
совершить, чтобы переместить единичный положительный заряд из точки A в
точку B (рис. 36.2).

\pandocbounded{\includegraphics[keepaspectratio]{images/_page_273_Figure_15.jpeg}}

Рис. 36 2 Потенциал точечного заряда в двумерном случае.
\texttt{\textbackslash{}alpha} --- работа, необходимая для перемещения
единичного заряда из A в BСтік можна інтерпретувати як від'ємне джерело,
так що стік -q створює потенціал

\begin{verbatim}
u\left(r\right)=-\frac{q}{2\pi}\ln\frac{1}{r}.
\end{verbatim}

Этим завершается обсуждение потенциалов точечных зарядов. Теперь мы
готовы к решению неоднородного уравнения методом функции Грина.\#
Уравнение Пуассона в кругеЗаймемся решением очень важной задачи(36.2)

\begin{verbatim}
(Y \cup \Pi)
\end{verbatim}

\texttt{\textquotesingle{}u\_\{rr\}\ +\ \textbackslash{}frac\{1\}\{r\}\ u\_r\ +\ \textbackslash{}frac\{1\}\{r\^{}2\}\ u\_\{\textbackslash{}theta\textbackslash{}theta\}\ =\ f(r,\ \textbackslash{}theta),\ \textbackslash{}quad\ 0\ \textless{}\ r\ \textless{}\ 1,\ \textbackslash{}quad\ 0\ \textbackslash{}le\ \textbackslash{}theta\ \textbackslash{}le\ 2\textbackslash{}pi,\textquotesingle{}\ \textquotesingle{}(\textbackslash{}Gamma\ Y)\textquotesingle{}\ \textquotesingle{}u(1,\ \textbackslash{}theta)\ =\ 0,\ \textbackslash{}quad\ 0\ \textbackslash{}le\ \textbackslash{}theta\ \textbackslash{}\ le\ 2\textbackslash{}pi.\textquotesingle{}\ **М**етод\ функции\ Грина\ (метод\ функции\ источника)\ состоит\ из\ двух\ шагов.1.\ Помещаем\ в\ точку\ \textquotesingle{}(\textbackslash{}rho,\ \textbackslash{}varphi)\textquotesingle{}\ точечный\ единичный\ заряд\ и\ находим\ потенциал\ \textquotesingle{}G(r,\ \textbackslash{}theta,\ \textbackslash{}rho,\ \textbackslash{}varphi)\textquotesingle{}\ ,\ который\ создается\ этим\ зарядом\ и\ обращается\ в\ нуль\ на\ границе.2.\ Суммируем\ индивидуальные\ отклики\ \textquotesingle{}G(r,\ \textbackslash{}theta,\ \textbackslash{}rho,\ \textbackslash{}phi)\textquotesingle{}\ ,\ взвешенные\ правой\ частью\ (плотностью\ зарядов)\ \textquotesingle{}f(r,\ \textbackslash{}theta)\textquotesingle{}\ ,\ по\ всему\ кругу.\ В\ результате\ получаем\ искомое\ решение}

\begin{verbatim}
u\left(r,\;\theta\right) = \int_{0}^{2\pi} \int_{0}^{1} G\left(r,\;\theta,\;\rho,\;\varphi\right) f\left(\rho,\;\varphi\right) \rho \; d\rho \; d\varphi.
\end{verbatim}

Найдем теперь функцию источника для нашей задачи.\# Определение функции
источника
\texttt{G(r,\ \textbackslash{}theta,\ \textbackslash{}rho,\ \textbackslash{}phi)}
Замінимо праву частину \texttt{f(r,\textbackslash{}theta)} рівняння
точковим джерелом величини 1, розташованим у деякій довольно обраній
точці \texttt{(\textbackslash{}rho,\ \textbackslash{}varphi)} .
Математично точкове джерело представлене дельта-функцією
\texttt{\textbackslash{}delta(r-\textbackslash{}rho,\ \textbackslash{}theta-\textbackslash{}varphi)\^{}1}
). Ми вважаємо, що дельта-функція дорівнює нулю всюди, крім точки
\texttt{(\textbackslash{}rho,\ \textbackslash{}varphi)} , де
розташований одиничний заряд. У мові сил дельта-функцію можна
інтерпретувати як точкова одиниця сили, прикладена до точки
\texttt{(\textbackslash{}rho,\ \textbackslash{}varphi)} . Наше завдання
--- знайти потенціал точкового заряду, якщо відомо, що він дорівнює нулю
на межі. Функція, яка задовольняє ці умови, називається функцією Гріна
або функцією джерела. Вона дорівнює \emph{відповіді} системи в точці
\texttt{(r,\ \textbackslash{}theta)} на \emph{джерело}, розташоване в
точці \texttt{(\textbackslash{}rho,\ \textbackslash{}varphi)} .
Складність розв'язання цієї проблеми полягає в тому, що бажаний
потенціал має зрівнятися з нулем на кордоні. Якби не ця умова, потенціал
можна було б легко знайти, адже ми вже знаємо, щоСм.: Владимиров В. С.
Обобщенные функции в математической физике --- М.: Наука, 1976. --- Прим
ред.потенциал точечного заряда равен

\begin{verbatim}
\frac{1}{2\pi} \ln \left( \frac{1}{r} \right)
\end{verbatim}

,где r --- расстояние от заряда до точки наблюдения.Искомый потенциал
может иметь следующий физический смысл:1. В теории теплопроводности ---
это равновесная температура внутри круга, если тепловой источник
размещается в точке (о. ф),а на границе поддерживается ну-левая
температура.- 2. Він описує відхилення мембрани від положення рівноваги,
якщо відхилення на межі дорівнюють нулю, а в точці (p.~f) мембрана
зігнута до дуже великої висоти. - 3. В електростатиці це розподіл
потенціалу всередині кола, якщо точковий позитивний заряд розташований у
точці \texttt{(\textbackslash{}rho,\ \textbackslash{}phi)} , а коло
заземлене, а його потенціал підтримується на нулі.

\pandocbounded{\includegraphics[keepaspectratio]{images/_page_275_Picture_11.jpeg}}

Единичная окружность РИС. 36 3. Функция Гряна 6 (г. 0. р, ф) точечного
источника, расположенного в точке (р, ф).Теперь найдем функцию Грина;
она будет похожа на функцию, изображенную на рис. 36.3.\# Построение
решенияШАГ 1. Поскольку функция

\begin{verbatim}
\frac{1}{2\pi} \ln \frac{1}{R}
\end{verbatim}

является потенциалом поля в точке
\texttt{p\ =\ (r,\ \textbackslash{}varphi)\ \textbackslash{}theta)} ,
которое создано единичным точечным зарядом, помещенным в точку
\texttt{Q\ =\ (\textbackslash{}rho,\ \textbackslash{}phi)} ,
\texttt{(R-paccmonnue\textbackslash{}\ between\textbackslash{}\ points\textbackslash{}\ P\textbackslash{}\ and\textbackslash{}\ Q)}
единственное, что нам осталось сделать---так изменить эту функцию, чтобы
она обратилась в нуль на границе.ЩАГ 2. Из эксперимента известно, что
эквипотенциальные линии точечных зарядов (положительных или
отрицательных) представляют собой окружности (рис. 36.4).Основна ідея
побудови функції Гріна така: потрібно розмістити інший (від'ємний) заряд
поза колом так, щоб загальний потенціал на колі r = 1 був сталим. Для
гему можна відняти цю константу з потенціалу і отримати нульовий
потенціал на межі. Очевидно. Таким чином знайдений потенціал матиме всі
властивості функції
\texttt{G(r,\ \textbackslash{}theta,\ \textbackslash{}rho,\ \textbackslash{}phi)}
. Залишається незрозумілим, у якій точці за межами кола слід розмістити
негативний заряд, щоб потенціал на колі став сталим. Можливо порівняльно
P И С

\pandocbounded{\includegraphics[keepaspectratio]{images/_page_276_Picture_3.jpeg}}

. 36 4. Потенциал поля, образованного двумя противоположно заряженными
частицами; a, b--- линии (окружности) постоянного потенциала.

\pandocbounded{\includegraphics[keepaspectratio]{images/_page_276_Picture_5.jpeg}}

\textbf{РИС.} 36.5. Заряды Q в \texttt{\textbackslash{}overline\{Q\}}
дают постоянный потенциал на окружности r=1.Легко побачити, що якщо
негативний заряд розміщений у точці
\texttt{\textbackslash{}bar\{Q\}\ =\ (\textbackslash{}bar\{\textbackslash{}rho\},\ \textbackslash{}bar\{\textbackslash{}phi\})\ =\ (1/\textbackslash{}rho,\ \textbackslash{}phi)}
, то повний потенціал

\begin{verbatim}
u(r, \theta) = \frac{1}{2\pi} \ln 1/R - \frac{1}{2\pi} \ln 1/\overline{R}
\end{verbatim}

буде сталою на колі r=1 Тут відстані від зарядів до точки (g, c)
позначаються R і \texttt{\textbackslash{}overline\{R\}} (див. рис.
36.5). Можна показати, що величина потенціалу на колі r = 1 визначається
виразом

\begin{verbatim}
-\frac{1}{2\pi}\ln\rho
\end{verbatim}

(положительная константа).С учетом этого результата строим функцию
Грина(36.3)

\begin{verbatim}
G(r, \theta, \rho, \phi) = \frac{1}{2\pi} \ln 1/R - \frac{1}{2\pi} \ln 1/\overline{R} + \frac{1}{2\pi} \ln \rho,
\end{verbatim}

Потенциал положентельного заряда, расположенного в тенциал границы
постоянного в тенциал границы постоянного в тенциал границы постоянного
в тенциал границы постоянного в тенциал границы постоянного в тенциал
границы постоянного в тенциал границы постоянного в тенциал границы
постоянного в тенциал границы постоянного в тенциал границы постоянного
в тенциал границы постоянного в тенциал границы постоянного в тенциал
границы постоянного в тенциал границы постоянного в тенциал границы
постоянного в тенциал границы постоянного в тенциал границы постоянного
в тенциал границы постоянного в тенциал границы постоянного в тенциал
границы постоянного в тенциал границы постоянного в тенциал границы
постоянного в тенциал границы постоянного в тенциал границы постоянного
в тенциал границы постоянного в тенциал границы постоянного в тенциал
границы постоянного в тенциал границы постоянного в тенциал границы
постоянного в тенциал границы постоянного в тенциал границы постоянного
в тенциал границы постоянного в тенциал границы постоянного в тенциал
границы постоянного в тенциал границы постоянного в тенциал границы
постоянного в тенциал границы постоянного в тенциал границы постоянного
в тенциал границы постоянного в тенциал границы постоянного в тенциал
границы постоянного в тенциал границы постоянного в тенциал границы
постоянного в тенциал границы постоянного в тенциал границы постоянного
в тенциал границы постоянного в тенциал границы постоянного в тенциал
границы постоянного в тенциал границы постоянного в тенциал границы
постоянного в тенциал границы постоянного в тенциал границы постоянного
в тенциал границы постоянного в тенциал границы постоянного в тенциал
границы постоянного в тенциал границы постоянного в тенциал границы
постоянного в тенциал границы постоянного в тенциал границы постоянного
в тенциал границы постоянного в тенциал границы постоянного в тенциал
границы постоян границы постоян границы постоян границы постоян границы
постоянгде

\begin{verbatim}
R = V \frac{r^2 - 2\rho r \cos(\theta - \phi) + \rho^2}{R},
\overline{R} = V \frac{r^2 - 2\frac{r}{\rho}\cos(\theta - \phi) + 1/\rho^2}{r^2 - 2\frac{r}{\rho}\cos(\theta - \phi) + 1/\rho^2}
\end{verbatim}

Формули, що виражають відстань між двома точками в полярних координатах.
Щоб знайти розв'язок початкової проблеми, потрібно накласти всі
імпульсні характеристики системи. Отже, переходимо до останнього
кроку.КРОК 3. Суперпозиція імпульсних характеристик Цей крок досить
простий. З відношення

\begin{verbatim}
u(r, \theta) = \int_{0}^{2\pi} \int_{0}^{1} G(r, \theta, \rho, \varphi) f(\rho, \varphi) \rho d\rho d\varphi
\end{verbatim}

получаем(36.4)

\begin{verbatim}
u(r, \theta) = \frac{1}{2\pi} \int_{0}^{2\pi} \int_{0}^{1} \ln (\rho \overline{R}/R) f(\rho, \varphi) \rho d\rho d\varphi.
\end{verbatim}

\begin{verbatim}
Формула (36.4) виражає розв'язок задачі Діріхле для рівняння Пуассона всередині кола через функцію Гріна цієї задачі. Якщо щільність заряду  f(r, \theta)  відома, то інтеграл можна обчислити, наприклад, чисельно.
\end{verbatim}

\section{ЗАУВАЖЕННЯ}\label{ux437ux430ux443ux432ux430ux436ux435ux43dux43dux44f}

\begin{enumerate}
\def\labelenumi{\arabic{enumi}.}
\tightlist
\item
  Задачу
\end{enumerate}

\begin{verbatim}
\begin{array}{llll} (\text{УЧ\Pi}) & \Delta u = 0, & 0 < r < 1, & 0 \leqslant \theta < 2\pi, \\ (\text{ГУ}) & u(1, \theta) = g(\theta), & 0 \leqslant \theta \leqslant 2\pi, \end{array}
\end{verbatim}

також можна розв'язати за допомогою методу \emph{функції Гріна}. У цьому
випадку рішення фіксується у формі

\begin{verbatim}
u(r, \theta) = \int_{0}^{2\pi} \frac{\partial G}{\partial r} (r, \theta, 1, \varphi) g(\varphi) d\varphi.
\end{verbatim}

\begin{verbatim}
Останнє відношення можна надати більш обчислювально корисну форму, явно виставляючи вираз у  \partial G/\partial r  . Внаслідок цього ми отримуємо інтегральну формулу Пуассона, знайдену раніше в лекції 33.(36.5)
\end{verbatim}

\begin{verbatim}
u(r, \theta) = \frac{1}{2\pi} \int_{0}^{2\pi} \left[ \frac{1 - r^2}{1 - 2r\cos(\theta - \varphi) + r^2} \right] g(\varphi) d\varphi.
\end{verbatim}

\begin{enumerate}
\def\labelenumi{\arabic{enumi}.}
\setcounter{enumi}{1}
\tightlist
\item
  Решение общей задачи Дирихле
\end{enumerate}

\begin{verbatim}
\begin{array}{llll} (\text{УЧП}) & \Delta u = f(r, \theta), & 0 < r < 1, & 0 \le \theta < 2\pi, \\ (\text{ГУ}) & u(1, \theta) = g(\theta), & 0 < \theta \le 2\pi, \end{array}
\end{verbatim}

находится как сумма решений (36.4) и (36.5).3. Метод функцій Гріна
дозволяє отримати розв'язки багатьох задач у областях різних форм. Однак
для кожної області і (точніше, для кожного оператора зліва від граничної
умови) для кожного рівняння потрібно знайти власну функцію Гріна, і це
не завжди легко зробити.4. Щоб дійсно знайти розв'язок за формулою
(36.4), у більшості випадків необхідно обчислити інтеграли на
комп'ютері.

\begin{verbatim}
#TASKS: Знайти потенціал точкового джерела у тривимірному просторі.2. Знайти функцію Гріна  G(x, y, \xi, \eta)  для рівняння Лапласа у верхній півплощині y > 0. Іншими словами, знайти потенціал у (x, y) у верхній півплощині, якщо заряд розташований у  (\xi, \eta)  , а потенціал на межі y=0 дорівнює нулю (див. рисунок).
\end{verbatim}

\pandocbounded{\includegraphics[keepaspectratio]{images/_page_278_Figure_11.jpeg}}

\begin{verbatim}
ПРИМІТКА. Якщо від'ємний заряд розмістити в точці  \overline{Q} = (\xi,  -\eta)  , очевидно, що негативний потенціал на прямій y = 0 буде нульовим. Те саме те, що функція Усмішки залежить від потенціалу цих двох зарядів.- 3. Використовуючи результати задачі 2, знайдемо розв'язок рівняння Пуассона  \Delta u = -k  у верхній півплощині з умовою нульової межі.
\end{verbatim}

\begin{itemize}
\tightlist
\item
  \begin{enumerate}
  \def\labelenumi{\arabic{enumi}.}
  \setcounter{enumi}{3}
  \tightlist
  \item
    Як би ви побудували функцію Зеленого для першого квадранта
    x\textgreater0,~y\textgreater0 ?
  \end{enumerate}
\item
  \begin{enumerate}
  \def\labelenumi{\arabic{enumi}.}
  \setcounter{enumi}{4}
  \tightlist
  \item
    Інший підхід до розв'язання рівняння Пуассона такий. Припустимо, ви
    хочете розв'язати проблему(УЧП)
  \end{enumerate}
\end{itemize}

\begin{verbatim}
\Delta u = 1
\end{verbatim}

\begin{verbatim}
,  0 < r < 1  ,  0 \le \theta < 2\pi  ,  (\Gamma Y)   u(1, \theta) = \sin \theta  ,  0 \le \theta \le 2\pi  .Сначала попытайтесь найти частное решение уравнения с помощью подстановки
\end{verbatim}

\begin{verbatim}
u_p = (r, \theta) = Ar^2
\end{verbatim}

\begin{verbatim}
.Потім підставляємо цей вираз у рівняння Пуассона і визначаємо константу  A_{\cdot}  числа  u_p  , що знаходиться в частині розв'язку, необхідно пожартувати поза розв'язком з огляду на  u=w+u_p  . Яка межа розуму, яку можна задовольнити функцією w? Знайте  w(r, \theta)  . Як вирішити проблему  u(r, \theta)  ? Зачекай. Retelno proanalizite vidpovid'. Зробіть перетин пеніса шкіри.# Частина 5. ЧИСЕЛЬНІ ТА ПРИБЛИЗНІ МЕТОДИЛекция 37# ЧИСЕЛЬНІ РОЗВ'ЯЗКИ (ЕЛІПТИЧНІ ЗАДАЧІ)МЕТА ЛЕКЦІЇ: Показати, як диференціальні рівняння в похідних похідних можна звести до систем алгебраїчних рівнянь, якщо похідні з частковими похідними замінити на наближення з скінченними різницями. Розв'язання цієї системи алгебраїчних рівнянь ітеративним методом дозволяє отримати наближене розв'язокНадається інформація про пакет ELLPACK, розроблений для розв'язання еліптичних приладівзадач общего вида на ЭВМ.На сьогодні ми ознайомилися з кількома методами розв'язання лінійних диференціальних рівнянь у похідних похідних. Однак більшість рівнянь, з якими ми працювали, були дуже простими. Форми регіонів, у яких розв'язувалися проблеми, також були простими. Однак багато завдань не можна спростити настільки, щоб їх зводити до певного набору шаблонних завдань. Такі задачі слід розв'язувати приблизно чисельними методами. Прогрес у сфері високопродуктивних комп'ютерів призвів до створення нових чисельних методів. Сьогодні такі нелінійні проблеми гідродинаміки, теорії пружності та теорії потенціалу вже розв'язані, про які ніхто навіть не думав десять років тому.Під загальною назвою «чисельні методи» поєднуються кілька різних підходів до розв'язання задач. Детальний опис усіх цих підходів можна знайти у книзі [4] з рекомендованого списку джерел. У цій та наступних двох лекціях ми покажемо, як розв'язувати еліптичні, гіперболічні та параболічні рівняння методом скінченної різниці.Спершу ознайомимося з поняттям скінченних різниць, а потім покажемо, як їх використовувати для розв'язання задачі Діріхле в квадраті.# Апроксимації скінченних різницьВспомним ряд Тейлора для функции f(x):
\end{verbatim}

\begin{verbatim}
f(x+h) = f(x) + f'(x)h + \frac{f''(x)}{2!}h^2 + \dots
\end{verbatim}

Якщо ми розірвемо цю суперечку на другому терміні, отримаємо

\begin{verbatim}
f(x+h) \cong f(x) + f'(x) h
\end{verbatim}

.Откуда

\begin{verbatim}
(37.1) f'(x) \cong \frac{f(x+h)-f(x)}{h}.
\end{verbatim}

\begin{verbatim}
Вираз праворуч називається похідною правої різниці. Вона апроксимує першу похідну f'(x) у точці x.В разложении Тейлора функции  f(\mathbf{r})  можно заменить h на -h и получить левую разностную производную(37.2)
\end{verbatim}

\begin{verbatim}
f'(x) \cong \frac{f(x) - f(x - h)}{h}
\end{verbatim}

\begin{verbatim}
.Buyutan  f(x - h) \cong f(x) - f'(x) h  \nif  f(x + h) \cong f( x) + f'(x) h  ,получаем центральную разностную производную(37.3)
\end{verbatim}

\begin{verbatim}
f'(x) \simeq \frac{1}{2h} [f(x+h) - f(x-h)].
\end{verbatim}

Якщо в ряду Тейлора залишиться ще один член, то центральну різницю
похідну для наближення f'\,'(x) можна отримати точно таким самим
способом(37.4)

\begin{verbatim}
f''(x) \cong \frac{1}{h^2} [f(x+h) - 2f(x) + f(x-h)].
\end{verbatim}

Тепер ми можемо розширити поняття наближення скінченних різниць до
\emph{часткових похідних}. Якщо виходити з розкладу Тейлора функції двох
змінних

\begin{verbatim}
u(x+h, y) = u(x, y) + u_x(x, y)h + u_{xx}(x, y)\frac{h^2}{2!} + \dots,
u(x-h, y) = u(x, y) - u_x(x, y)h + u_{xx}(x, y)\frac{h^2}{2!} - \dots,
\end{verbatim}

можно получить следующие аппроксимации частных производных:

\begin{verbatim}
\begin{split} u_x(x,\ y) & \cong \frac{u(x+h,\ y)-u(x,\ y)}{h}, \\ u_{xx}(x,\ y) & \cong \frac{1}{h^2} \big[ u(x+h,\ y)-2u(x,\ y)+u(x-h,y) \big], \\ u_u(x,\ y) & \cong \frac{u(x,\ y+k)-u(x,\ y)}{k}, \\ u_{uy}(x,\ y) & \cong \frac{1}{h^2} \big[ u(x,\ y+k)-2u(x,\ y)+u(x,\ y-k) \big]. \end{split}
\end{verbatim}

В этих формулах частные производные аппроксимируются правыми,
центральными и левыми разностными производными, но в данной лекции мы
будем пользоваться только центральными аппроксимациями.

\pandocbounded{\includegraphics[keepaspectratio]{images/_page_282_Figure_5.jpeg}}

Рис. 37.1 Сетка для задачи Дирихле внутри квадрата.Чтобы познакомиться с
основными правилами использования конечно-разностных аппроксимаций,
рассмогрим простую задачу Дирихле.Решение задачи Дирихле методом
конечных разностейПусть необходимо решить задачу Дирихле(ЗЧП)

\begin{verbatim}
u_{xx} + u_{yy} = 0
\end{verbatim}

\texttt{,\ \textquotesingle{}0\ \textless{}\ x\ \textless{}\ 1\textquotesingle{}\ ,\ \textquotesingle{}0\ \textless{}\ y\ \textless{}\ 1\textquotesingle{}\ ,\ \textquotesingle{}u\ =\ 0\textquotesingle{}\ ,\ на\ верхней\ и\ боковых\ сторонах\ квадрата,\ \textquotesingle{}u(x,\ 0)\ =\ \textbackslash{}sin(x\textbackslash{}pi)\textquotesingle{}\ ,\ \textquotesingle{}0\ \textbackslash{}le\ \textbackslash{}le\ x\ \textbackslash{}le\ 1\textquotesingle{}\ ,\ \textquotesingle{}y\ =\ 0\textquotesingle{}\ .Построим\ в\ плоскости\ (x,\ y)\ сетку,\ как\ показано\ на\ рис.\ 37.1.\ Удобно\ (особенно,\ если\ мы\ собираемся\ применять\ ЭВМ)непользовать\ следующие\ обозначения:}

\begin{verbatim}
\begin{split} u\left(x,\,y\right) &= u_{ij}\,,\\ u\left(x,\,y+k\right) &= u_{i+1,j}\,,\\ u\left(x,\,y-k\right) &= u_{i-1,j}\,,\\ u\left(x-h,\,y\right) &= u_{i,j-1}\,,\\ u\left(x+h,\,y\right) &= u_{i,j+1}\,,\\ u_{x}\left(x,\,y\right) &= \frac{1}{2h}\left(u_{i,j+1}-u_{i,j-1}\right)\,,\\ u_{y}\left(x,\,y\right) &= \frac{1}{2k}\left(u_{i+1,j}-u_{i-1,j}\right)\,,\\ u_{xx}\left(x,\,y\right) &= \frac{1}{h^{2}}\left(u_{i,j+1}-2u_{i,j}+u_{i,j-1}\right)\,,\\ u_{yy}\left(x,\,y\right) &= \frac{1}{h^{2}}\left(u_{i+1,j}-2u_{i,j}+u_{i-1,j}\right)\,. \end{split}
\end{verbatim}

Основная идея решения задачи Дирихле основана на замене частных
производных в уравнении Лапласа

\begin{verbatim}
u_{xx} + u_{yy} = 0
\end{verbatim}

их конечно-разностными аппроксимациями. Проделав это и воспользовавшись
компактными обозначениями, получаем следующее разностное уравнение:

\begin{verbatim}
\Delta u = \frac{1}{h^2} (u_{i,j+1} - 2u_{i,j} + u_{i,j-1}) + \frac{1}{k^2} (u_{i+1,j} - 2u_{i,j} + u_{i-1,j}) = 0.
\end{verbatim}

В случае когда h=k, уравнение Лапласа приводится к виду (37.6)
\texttt{(u\_\{i+1,j\}+u\_\{i-1,j\}+u\_\{i,j+1\}+u\_\{i,j+1\}-4u\_\{ij\})=0.}
Разрешив его относительно \texttt{u\_\{II\}} , получим

\begin{verbatim}
u_{i,j} = \frac{1}{4} (u_{i+1,j} + u_{i-1,j} + u_{i,j+1} + u_{i,j-1}).
\end{verbatim}

Зверніть увагу, що в цьому співвідношенні всі значення
\texttt{u\_\{tj\}} беруть для внутрішніх вузлів сітки. Згідно з останнім
співвідношенням, розв'язок \texttt{u\_\{ij\}} апроксимується середнім
значенням розв'язку у чотирьох сусідніх точках. Відповідно, ми можемо
розробити чисельний метод для розв'язання цієї задачі.Алгоритм
численного решения задачи Дирихле (метод Либманна)ШАГ 1. Присвоим
величинам \texttt{u\_\{ij\}} во внутренних узлах численное вначение,
равное среднему всех значений граничных условий.КРОК 2. Ми перерахуємо
значення у всіх внутрішніх точках сітки, замінивши старе значення на
середнє значення чотирьох суміжних точок. Не дуже важливо, як буде
організовано підрахунок процесу, але зазвичай він здійснюється рядками
(або колонками). Після кількох операцій процес збігається до приблизного
розв'язання задачі. Швидкість збіжності цього алгоритму низька, але її
можна збільшити кількома способами. Рекомендуємо зацікавленим читачам
звернутися до книги {[}4{]} зі списку рекомендованої літератури. \#
ЗАУВАЖЕННЯ 1. Если взять сетку с чегырьмя внутренними точками (m=n=4),
то система (37.6) запишется в виде

\begin{verbatim}
-4u_{22} + 0 + \sin(\pi/3) + u_{23} + u_{32} = 0,
-4u_{23} + u_{22} + \sin(2\pi/3) + 0 + u_{33} = 0,
-4u_{32} + 0 + u_{22} + u_{33} + 0 = 0,
-4u_{33} + u_{32} + u_{23} + 0 + 0 = 0.
\end{verbatim}

Це потрібно вирішувати у зв'язку з \texttt{u\_\{22\}} ,
\texttt{u\_\{21\}} , \texttt{u\_\{32\}} та \texttt{u\_\{33\}} .
Розв'язок такої системи можна знайти за допомогою \emph{ітеративного
методу}. Метод Лібмана також є одним із ітеративних методів.- 2. Якщо
зменшити кроки сітки на h і k (кількість вузлів сітки збільшиться), то
отримаємо систему рівнянь типу (37,7), але більших розмірів.

\begin{itemize}
\tightlist
\item
  \begin{enumerate}
  \def\labelenumi{\arabic{enumi}.}
  \setcounter{enumi}{2}
  \tightlist
  \item
    Систему рівнянь (37.7) можна переписати у матричній формі
  \end{enumerate}
\end{itemize}

\begin{verbatim}
\begin{bmatrix} -4 & 1 & 1 \\ 1 & -4 & 0 & 1 \\ 1 & 0 & -4 & 1 \\ 0 & 1 & 1 & -4 \end{bmatrix}
\begin{bmatrix} u_{22} \\ u_{23} \\ u_{32} \\ u_{33} \end{bmatrix}
=
\begin{bmatrix} -\sin{(\pi/3)} \\ -\sin{(2\pi/3)} \\ 0 \\ 0 \end{bmatrix}
=
\begin{bmatrix} -0.86 \\ -0.86 \\ 0 \\ 0 \end{bmatrix}
\end{verbatim}

.

Зі збільшенням кількості рівнянь (наприклад, до 1000) матриця
коефіцієнтів стає розрідженою, тобто містить багато нулів. Системи з
розрідженими матрицями розв'язуються спеціальними методами. Зазвичай
використовуються ітеративні методи, такі як метод Якобі, метод
Гаусса-Зайделя або послідовний метод верхньої релаксації.4. При
розв'язанні задачі Неймана ті \emph{похідні}, що входять до граничної
умови, також слід замінити різницею апроксимаціями.5. Таким чином ви
зможете вирішувати проблемиa)

\begin{verbatim}
u_{xx} + u_{yy} = f(x, y)
\end{verbatim}

\begin{enumerate}
\def\labelenumi{\arabic{enumi})}
\setcounter{enumi}{5}
\tightlist
\item
\end{enumerate}

\begin{verbatim}
xu_{xx} + u_{yy} + 2u = \sin(x - y)
\end{verbatim}

\begin{enumerate}
\def\labelenumi{\Alph{enumi})}
\setcounter{enumi}{1}
\tightlist
\item
\end{enumerate}

\begin{verbatim}
\sin xu_{xx} + u_{xy} + 3u = 0
\end{verbatim}

(гетерогенні рівняння з змінними коефіцієнтами), (рівняння зі змінними
коефіцієнтами).6. Якщо область, в якій розв'язується задача, має
неправильну форму, її можна покрити сіткою і наближити розв'язок у
точках, найближчих до межі, інтерполюючи граничні умови. Після цього
задачу вирішують звичайним способом (див. рисунок 37.2).

\pandocbounded{\includegraphics[keepaspectratio]{images/_page_285_Figure_8.jpeg}}

□ новые ГУ находятся интерполяцией ГУ на кривой Рис. 37 2.- 7. Деякі
журнали публікують тексти програм для розв'язання диференціальних
рівнянь у похідних на комп'ютерах. Нижче наведено короткий список цих
журналів:

\begin{itemize}
\tightlist
\item
  \begin{enumerate}
  \def\labelenumi{\alph{enumi})}
  \tightlist
  \item
    Транзакції ACM на математичному програмному забезпеченні,
  \end{enumerate}
\item
  \begin{enumerate}
  \def\labelenumi{\arabic{enumi})}
  \setcounter{enumi}{5}
  \tightlist
  \item
    Комп'ютерний журнал,
  \end{enumerate}
\item
  \begin{enumerate}
  \def\labelenumi{\Alph{enumi})}
  \setcounter{enumi}{1}
  \tightlist
  \item
    Numerische Mathematik,
  \end{enumerate}
\item
  \begin{enumerate}
  \def\labelenumi{\alph{enumi})}
  \setcounter{enumi}{17}
  \tightlist
  \item
    БІТКрім того, нещодавно було розроблено пакет застосунків ELLPACK
    для розв'язання досить широкого класу задач на лінії прикордонних
    ліній. Цей пакет дозволить вам розв'язувати різноманітні дво- та
    тривимірні задачі урізні координатні системи для довільних меж і
    спільних граничних умов. Користувачу пропонується широкий спектр
    різних методів розв'язання 1).
  \end{enumerate}
\end{itemize}

\section{\texorpdfstring{ЗАДАЧІ - 1 Отримаємо наближення (39,4) для
другої похідної f'\,`(x) f'\,'(x) = \{\texttt{f(x+h)} 2f(x) +
\texttt{f(x-h)}\}/h\^{}2
.}{ЗАДАЧІ - 1 Отримаємо наближення (39,4) для другої похідної f'\,`(x) f'\,'(x) = \{f(x+h) 2f(x) + f(x-h)\}/h\^{}2 .}}\label{ux437ux430ux434ux430ux447ux456---1-ux43eux442ux440ux438ux43cux430ux454ux43cux43e-ux43dux430ux431ux43bux438ux436ux435ux43dux43dux44f-394-ux434ux43bux44f-ux434ux440ux443ux433ux43eux457-ux43fux43eux445ux456ux434ux43dux43eux457-fx-fx-fxh-2fx-fx-hh2-.}

\begin{itemize}
\tightlist
\item
  \begin{enumerate}
  \def\labelenumi{\arabic{enumi}.}
  \setcounter{enumi}{1}
  \tightlist
  \item
    Виконайте дві ітерації задачі Діріхле (37.5) за ітеративним методом
    Лібмана.
  \end{enumerate}
\item
  \begin{enumerate}
  \def\labelenumi{\arabic{enumi}.}
  \setcounter{enumi}{2}
  \tightlist
  \item
    До якої алгебранічної системи буде зведена задача Діріхле для
    рівнянь Пуассона в квадраті
  \end{enumerate}
\end{itemize}

\texttt{(УЧ\textbackslash{}Pi)}

\begin{verbatim}
u_{xx} + u_{yy} = f(x,y), \quad 0 < x < 1, \quad 0 < y < 1, (\Gamma Y) u(x,y) = g(x,y) на границе, если ее решать ме-Скінченні відмінності? 4. Розв'язати задачу 3, якщо(УЧП)
\end{verbatim}

\begin{verbatim}
u_{xx} + u_{yy} + 2u = 0
\end{verbatim}

, 0 \textless{} x \textless{} 1 , 0 \textless{} y \textless{} 1 , (ГУ)
\texttt{u(x,\ y)} = \texttt{g(x,\ y)} на границе.5. Как бы вы решали
задачу Неймана внутри квадрата(УЧП)

\begin{verbatim}
u_{xx} + u_{yy} = 0
\end{verbatim}

, 0 \textless{} x \textless{} 1 , 0 \textless{} y \textless{} 1 , на
верхней, нижней и левой стороне квадрата,

\begin{verbatim}
\begin{cases} u = 0 & \text{квадрата,} \\ \frac{\partial u}{\partial x}(1, y) = 1, & 0 \leqslant y \leqslant 1, \end{cases}
\end{verbatim}

методом конечных разностей?6. Постройте блок-схему алеоритма решения
задачи Дирихле в квадрате

\begin{verbatim}
(\text{УЧП})
\end{verbatim}

\begin{verbatim}
u_{xx} + u_{yy} = f(x, y), \quad 0 < x < 1, \quad 0 < y < 1, (\text{ГУ}) u(x, y) = g(x, y) на границе,- з довільною кількістю вузлів сітки. Якщо ви знаєте мову програмування, напишіть програму для виконання цих обчислень.
\end{verbatim}

\begin{itemize}
\tightlist
\item
  \begin{enumerate}
  \def\labelenumi{\arabic{enumi})}
  \tightlist
  \item
    Російською мовою публікуються статті з чисельних методів розв'язання
    диференціальних рівнянь у численних журналах і збірках. Наприклад:
    «Журнал обчислювальної математики та математичної фізики»,
    видавництво «Наука», «Журнал інженерної фізики», видавництво
    Академії наук БРСР, збірки серії «Обчислювальні методи і
    програмування», Видавництво Московського державного університету.
    Багато публікацій, переважно статей, присвячені наборам прикладних
    проблем. Як приклад, можна вказати: Горбунов-Посадов М. М., Карпов
    В. Я, Корягін та ін. Пакет SAFRA. Програмне забезпечення для
    обчислювальних експериментів. У: Пакети прикладних програм.
    Обчислювальний експеримент. Москва, Наука, 1983. Примітка.
    Переклад.\# МЕТАЛЕКЦІЯ ПРО ЯВНІ РІЗНИЦІ СХЕМ: Введіть поняття явної
    різниці та покажіть, як її можна використовувати для розв'язання
    гіперболічних і параболічних задач. Основна ідея полягає в тому, що
    після заміни рівняння типу
  \end{enumerate}
\end{itemize}

\begin{verbatim}
u_t = u_{xx}
\end{verbatim}

\begin{verbatim}
Її апроксимація скінченної різниці дає формули, які явно виражають значення розв'язку за один момент часу через значення розв'язку в попередній момент часу. Таким чином, змішану задачу для параболічного або гіперболічного рівняння можна розв'язати, послідовно обчислюючи розв'язки для всіх наступних точок часу.Явна схема не позбавлена недоліків. Якщо необхідно підвищити точність наближення похідних, то при зміні сегмента збільшується не лише обсяг обчислень, а й похибки округлення.Попередня лекція була присвячена розв'язанню еліптичних задач на межі лінії (стаціонарних задач). У таких задачах необхідно знайти розв'язок диференціального рівняння в похідних у заданій області простору, якщо розв'язок або його похідна задано на межі області визначення Приблизне розв'язання еліптичних задач зводиться до розв'язку системи алгебраїчних рівнянь для значень функції у внутрішніх вузлах сітки. Іншими словами, значення розв'язку у всіх внутрішніх вузлах визначаються одночасно.У цій лекції ми розглянемо диференціальні методи розв'язання задач, залежних від часу. Основна ідея полягає в тому, що якщо розв'язок відомий у початковий момент часу, то можливо знайти розв'язок за допомогою схеми підрахунку t=\Delta t,\ 2\Delta t,\ 3\Delta t,\ldots Замінивши часткові похідні часу та просторову змінну на скінченнорізні похідні, можна отримати явні вирази для u_{tj} через значення функції u у попередні моменти часу. Такий процес називається явною схемою підрахунку з тикуванням.Щоб показати цей метод у дії, розглянемо типову задачу теплопровідності.# Явная схема для уравнения теплопроводностиРозглянемо проблему теплопровідності у стрижні, початкова температура якого дорівнює нулю. Нехай температура лівого кінця фіксована, а на правому кінці відбувається теплообмін із навколишнім середовищем, щоб тепловий фланець був пропорційний P и с
\end{verbatim}

\pandocbounded{\includegraphics[keepaspectratio]{images/_page_288_Figure_3.jpeg}}

\begin{verbatim}
. 38.1. Сетка для уравнения теплопроводности.різниця температур між кінцем стрижня і середовищем. Нехай температура середовища задається функцією g\left(t\right) . Іншими словами, ми розв'язуємо проблемузадачу (УЧП)
\end{verbatim}

\begin{verbatim}
u_t = u_{xx}, \quad 0 < x < 1, \quad 0 < t < \infty,
\end{verbatim}

(38.1) (ГУ)

\begin{verbatim}
\begin{cases} u(0, t) = 1, \\ u_x(1, t) = -[u(1, t) - g(t)], \\ u(x, 0) = 0, \quad 0 \le x \le 1. \end{cases}
\end{verbatim}

Щоб розв'язати цю задачу, ми побудуємо прямокутну сітку методом
скінченної різниці, вершини якої визначаються формулами (див. рис. 38.1)

\begin{verbatim}
x_j = jh
\end{verbatim}

, j = 0, 1, 2, \ldots, n ,
\texttt{y\_i\ =\ ik\ ,\ i\ =\ 0,\ 1,\ 2,\ ...,\ m}.Зверніть увагу, що
значення \texttt{u\_\{ij\}} зліва і внизу сітки на рис. 38.1 відомі з
граничних і початкових умов, і наша вадача полягає у знаходженні решти
значень \texttt{u\_\{ij\}} . Щоб розв'язати цю задачу, ми замінимо
часткові похідні рівняння теплопровідності на їхні наближення з
скінченними різницями

\begin{verbatim}
u_{t} = \frac{1}{k} [u(x, t+k) - u(x, t)] = \frac{1}{k} [u_{t+1,j} - u_{t,j}],
u_{xx} = \frac{1}{h^{2}} [u(x+h, t) - 2u(x, t) + u(x-h, t)] = \frac{1}{h^{2}} [u_{t,j+1} - 2_{t,j} + u_{t,j-1}].
\end{verbatim}

Подставим эти выражения в уравнение \texttt{u\_t\ =\ u\_\{xx\}} и
разрешим получившееся уравнение относительно значений функции на верхнем
временном слое. В результате получаем(38.2)

\begin{verbatim}
u_{t+1,j} = u_{t,j} + \frac{k}{h^2} \left[ u_{t,j+1} - 2u_{t,j} + u_{t,j-1} \right].
\end{verbatim}

Це та формула, яку ми шукаємо, оскільки вона виражає розв'язок у певний
момент часу через розв'язок у попередній момент часу (індекс \emph{i}
означає часову змінну). Фіг. 38.1 підкреслює значення, включені до цієї
формули.Тепер можна починати розрахунки. Однак спочатку потрібно
апроксимувати похідну в граничній умові праворуч

\begin{verbatim}
a_x(1, t) = -[u(1, t) - g(t)].
\end{verbatim}

В результате аппроксимации получаем(38.3)

\begin{verbatim}
\frac{1}{h}[u_{t,n} - u_{t,n-1}] = -[u_{t,n} - g_t].
\end{verbatim}

где значения \texttt{g\_i\ =\ g(ik)} известны. Здесь мы заменили
\texttt{u\_x(1,\ t)} левой разностной производной, поскольку правая
разностная производная потребовала бы значений функции за пределами
сетки. Из (38.3) находим(38.4)

\begin{verbatim}
u_{t,n} = \frac{u_{t,n-1} - hg_t}{1 + h}.
\end{verbatim}

\texttt{Формулы\ (38.2)\ и\ (38.4)\ позволяют\ начать\ вычисления.Алгоритм\ вычислений\ по\ явной\ схемеШАГ\ 1.\ Находим\ решение\ на\ сеточном\ слое\ \textquotesingle{}t\ =\ \textbackslash{}Delta\ t\textquotesingle{}\ ,\ используя\ явную\ формулу\ (см.\ рис.\ 38.2)}

\begin{verbatim}
u_{2,j} = u_{1,j} + \frac{k}{h^2} [u_{1,j+1} - 2u_{1,j} + u_{1,j-1}], \quad j = 2, 3, \ldots, n-1.
\end{verbatim}

ШАГ 2. Величину \texttt{u\_\{2,n\}} находим по формуле (38.4)

\begin{verbatim}
u_{2,n} = \frac{u_{2,n-1} + hg_2}{1+h}.
\end{verbatim}

\texttt{Після\ виконання\ кроків\ 1\ і\ 2\ ми\ отримуємо\ рішення\ для\ \textquotesingle{}t=\textbackslash{}Delta\ t\textquotesingle{}\ .\ Щоб\ отримати\ розв\textquotesingle{}язок\ у\ \textquotesingle{}t=2\textbackslash{}Delta\ t\textquotesingle{}\ (другий\ рядок\ знизу\ на\ рис.\ 38.2),\ повторіть\ кроки\ 1\ і\ 2,\ рухаючись\ на\ одну\ лінію\ вгору,\ тобто\ збільшуючи\ t\ на\ 1\ і\ використовуючи\ \textquotesingle{}u\_\{t,j\}\textquotesingle{}\ з\ попередньої\ лінії.\ Розв\textquotesingle{}язок\ обчислюється\ таким\ самим\ способом\ у\ наступні\ \textquotesingle{}t=3\textbackslash{}Delta\ t,\textbackslash{}\ 4\textbackslash{}Delta\ t\textbackslash{}dots\textquotesingle{}\ Для\ того\ чтобы\ помочь\ читателю\ в\ проведении\ расчетов\ по\ этому\ методу,\ на\ рис.\ 38.3\ приведена\ достаточно\ подробная\ блоксхема\ алгоритма.\ Те,\ кто\ не\ очень\ хорошо\ знаком\ с\ блок-схемами.}

\pandocbounded{\includegraphics[keepaspectratio]{images/_page_290_Figure_3.jpeg}}

Рис. 38.2. Иллюстрация явной разностной схемы: n--- число узлов сетки
вдоль оси \texttt{x;\textbackslash{}\ m} --- число узлов сетки вдоль оси
\texttt{y;\textbackslash{}\ h=1/(m-1),\textbackslash{}\ k=T/(m-1)} могут
счигать их промежуточным звеном между вычислительными алгоритмами и
подробными программами для ЭВМ. Блок-схемы дают строгое описание
процесса вычислений. \# ЗАУВАЖЕННЯ - Явна схема має серйозний недолік.
Якщо часовий крок достатньо великий порівняно з кроком x, похибки
округлення можуть стати настільки великими, що отриманий розв'язок
втрачає сенс. Відношення кроків у t і x залежить від рівняння та
граничних умов, але в загальному випадку крок у часі має бути значно
меншим за крок у координаті. У книзі {[}4{]} доведено, що для
застосування явної схеми необхідно виконати умову k/h2 ≤ 0,5.
\texttt{-\ 2.\ Існує\ таке\ правило:\ якщо\ зменшити\ кількість\ кроків\ \textquotesingle{}\textbackslash{}Delta\ t\textquotesingle{}\ і\ \textquotesingle{}\textbackslash{}Delta\ x\textquotesingle{}\ ,\ то\ похибка\ у\ наближенні\ часткових\ похідних\ за\ скінченними\ різницями\ також\ зменшиться,\ але\ чим\ менша\ сітка,\ тим\ більше\ обчислень\ потрібно\ виконати,\ і,\ відповідно,\ тим\ більшими\ будуть\ похибки\ округлення.\ Це\ явище\ ілюструється\ на\ рис.\ 38.5.}
- 3. Для гіперболічної задачі(УЧП)

\begin{verbatim}
u_{tt} = u_{xx}, \quad 0 < x < 1, \quad 0 < t < \infty,
\end{verbatim}

(ГУ)

\begin{verbatim}
\begin{cases} u(0, t) = g_1(t), \\ u(1, t) = g_2(t), \end{cases}
\quad 0 < t < \infty,
\end{verbatim}

(HY)

\begin{verbatim}
\begin{cases} u(x, 0) = \varphi(x), \\ u_t(x, 0) = \psi(x), \end{cases}
\quad 0 \leqslant x \leqslant 1,
\end{verbatim}

Також можливо побудувати явну різницю схеми. Для цього ми апроксимуємо
\texttt{u\_\{tt\}} і \texttt{u\_\{xy\}} центральних різницьких
пропозицій И С

\pandocbounded{\includegraphics[keepaspectratio]{images/_page_291_Figure_4.jpeg}}

. 38.3. Алгоритм решения задачи по явной схеме.водными

\begin{verbatim}
u_{tt} \simeq \frac{1}{h^2} [u(x, t+k) - 2u(x, t) + u(x, t-k)],
u_{xx} \simeq \frac{1}{h^2} [u(x+h, t) - 2u(x, t) + u(x-h, t)].
\end{verbatim}

а начальное условие--- по схеме

\begin{verbatim}
u_t(x, 0) \approx \frac{1}{k} [u(x, k) - u(x, 0)] = \frac{1}{k} [u(x, k) - \varphi(x)].
\end{verbatim}

В результате для вычисления величины \texttt{u(x,\ t+k)} получаем
следующую явную схему:(38.5)

\begin{verbatim}
u(x, t+k) = 2u(x, t) - u(x, t-k) + + (k/h)^2 [u(x+h, t) - 2u(x, t) + u(x-h, t)].
\end{verbatim}

З (38.5) видно, що для обчислення розв'язку на заданому часовому шарі
необхідно знати розв'язок на двох попередніх шарах. Отже, щоб почати
рахувати, потрібно використовувати P И С

\pandocbounded{\includegraphics[keepaspectratio]{images/_page_292_Figure_6.jpeg}}

. 38.4. Использование массивов при реализации явной схемы.начальным
условнем на скорость

\begin{verbatim}
\frac{1}{k}\left[u\left(x,\,k\right)-\varphi\left(x\right)\right]=\psi\left(x\right),
\end{verbatim}

\begin{verbatim}
з якої ми виводимо  u(x,  k) = \varphi(x) + k\psi(x)  , тобто цінність розв'язку на  t = \Delta t  . Розв'язок у наступні часи можна знайти за явною формулою (38.5).
\end{verbatim}

\section{ЗАДАЧІ}\label{ux437ux430ux434ux430ux447ux456}

\begin{enumerate}
\def\labelenumi{\arabic{enumi}.}
\tightlist
\item
  Постройте явную конечно-разностную схему для задачи(УЧП)
\end{enumerate}

\begin{verbatim}
u_t = u_{xx}
,  0 < x < 1  ,  0 < t < \infty  ,
\end{verbatim}

(ГУ)

\begin{verbatim}
\begin{cases} u(0, t) = 0, \\ u(1, t) = 0, \\ u(x, 0) = \sin(\pi x), \quad 0 \le x \le 1. \end{cases}
\end{verbatim}

\begin{verbatim}
(HУ) u(x, 0) = \sin(\pi x), \quad 0 \le x \le 1. Полагая h = \Delta x = 0.1 , найдите решение при t_1 = 0.005 , t_2 = 0.010 , t_3 = 0.015 . Постройте график полученного решения на сетке x = 0; 0.1; 0.2; 0.3; . .; 0.9: 1 при t = 0.015.
\end{verbatim}

\pandocbounded{\includegraphics[keepaspectratio]{images/_page_293_Figure_6.jpeg}}

Фіг. 38.5. Залежність повної похибки від кроку сітки.- 2. Розв'яжіть
задачу 1 аналітично (ділячи змінні). Знайдіть значення розв'язку у
вузлах сітки x=0; 0,1; 0,2; 0,3; \ldots; 0,9; 1 при t=0,015. Порівняйте
ці результати з результатами чисельного розв'язку задачі 1. (Ви можете
скористатися калькулятором або написати невелику комп'ютерну програму
для табуляції аналітичного розв'язку.)

\begin{itemize}
\tightlist
\item
  \begin{enumerate}
  \def\labelenumi{\arabic{enumi}.}
  \setcounter{enumi}{2}
  \tightlist
  \item
    Відповідно до примітки 3, побудуйте блок-схему для розв'язання
    гіперболічної задачі.
  \end{enumerate}
\item
  \begin{enumerate}
  \def\labelenumi{\arabic{enumi}.}
  \setcounter{enumi}{3}
  \tightlist
  \item
    Розв'яжіть задачу 1 з новою граничною умовою при x=1
    \texttt{u\_x}(1, t) = -{[}\texttt{u(1,\ t)} 1{]}. \# СХЕМИ НЕЯВНИХ
    РІЗНИЦЬ (СХЕМА КРАНКА--НІКОЛСОНА)
  \end{enumerate}
\end{itemize}

МЕТАЛЕКЦІЇ: Показати, як задачі, що містять функції, залежні від часу,
можна розв'язати за допомогою неявної різницької схеми. У цій схемі, як
і в попередній схемі, часткові похідні замінюються скінченно-різницею,
але \texttt{u\_\{l+1,\textbackslash{}\ J\}} більше не виражаються явно
через значення на попередніх шарах. Тепер, щоб визначити значення
\texttt{u\_\{l+1,\textbackslash{}\ J\}} , необхідно розв'язати систему
рівнянь. Іншими словами, щоб визначити розв'язок на кожному часовому
nage, необхідно розв'язати систему рівнянь.Перевага неявних схем над
явними полягає в тому, що в неявних схемах танг сітки можна зробити
достатньо великим без страху, що помилки округлення «знищать»
розв'язок.Для розв'язання параболічної задачі ми використаємо відому
схему \emph{Кранк}-\emph{Ніколсона}.Як уже згадувалося в попередній
лекції, щоб явна схема працювала надійно, часовий крок сітки має бути
невеликим. Зокрема, якщо просте завдання для рівняння
теплопровідності(39.1)

\begin{verbatim}
\begin{array}{ccccccccccccccccccccccccccccccccccc
\end{verbatim}

\begin{verbatim}
розв'язується за явною схемою, тоді кроки сітки \Delta t і \Delta x повинні задовольняти умову
\end{verbatim}

\begin{verbatim}
\frac{\Delta t}{(\Delta x)^2} \leqslant 0.5.
\end{verbatim}

\begin{verbatim}
Інакше метод буде чисельно нестабільним (а помилки округлення зростуть нескінченно). (Питання чисельної стабільності розглядаються у [53].) Іншими словами, якщо відстань сітки на осі x вибрана \Delta x=0,1 , то часовий крок \Delta t не може бути більшим за 0.5\Delta x^2=0.005 . Це означає, що ви можете пройти від t=0 до t=1, зробивши 200 кроків.Однак існують такі схеми (неявні схеми), які дозволяють виконувати обчислення з набагато більшим кроком, тоді якПравда, обсяг розрахунків, які виконуються на одному кроці, зростає. Такі схеми можна використовувати для виконання обчислень із досить великим кроком, але на кожному кроці потрібно розв'язати систему алгебраїмних рівнянь. Щоб проілюструвати цей метод, розв'яжемо наступну проблему теплопровідності.# Неявная схема для уравнения теплопроводностиРассмотрим задачу(39 2)
\end{verbatim}

\begin{verbatim}
\begin{array}{ccccccccccccccccccccccccccccccccccc
\end{verbatim}

Використаємо наступні апроксимації скінченних різниць для часткових
похідних \texttt{u\_t} та \texttt{u\_\{xx\}} :

\begin{verbatim}
\begin{split} u_t(x,\ t) &= \frac{1}{k} \big[ u(x,\ t+k) - u(x,\ t) \big], \\ u_{xx}(x,\ t) &= \frac{\lambda}{h^2} \big[ u(x+h,\ t+k) - 2u(x,\ t+k) + u(x-h,\ t+k) \big] + \\ &\quad + \frac{(1-\lambda)}{h^2} \big[ u(x+h,\ t) - 2u(x,\ t) + u(x-h,\ t) \big], \end{split}
\end{verbatim}

\begin{verbatim}
де \lambda взято з сегмента [0,1]. Зверніть увагу, що u_{xx} наближається зваженим середнім центральних різницевих похідних у момент часу t і t+k. \lambda=0.5 отримується звичайне середнє для цих двох центральних похідних, і при \lambda=0.75 одна з різницевих похідних бере вагу 0,75, а інша — 0,25. У \lambda=0 отримується звичайна явна схема, яка була обговорена в попередній лекції.Після заміни часткових похідних u_t і u_{xx} у задачі (39.2) отримуємо різницю (див. рисунок 39.1).Разностное уравнение
\end{verbatim}

\begin{verbatim}
\frac{1}{k} (u_{i+1, j} - u_{i, j}) = \frac{\lambda}{h^2} (u_{i+1, j+1} - 2u_{i+1, j} + u_{i+i, j-1}) + \frac{(1-\lambda)}{h^2} (u_{i, j+1} - 2u_{i, j} + u_{i, j-1}),
\end{verbatim}

(39.3)

\begin{verbatim}
(\Gamma Y) \begin{cases} u_{i, 1} = 0, & i = 1, 2, ..., m, \\ u_{i, n} = 0, & i = 1, 2, ..., m, \\ u_{i, j} = 1, & j = 2, ..., n-1. \end{cases}
\end{verbatim}

298Перенесіть усі невідомі значення u з верхнього часового шару (i+1) на
ліву частину рівняння (39.3) і отримайте (39.4)

\begin{verbatim}
-\lambda r u_{i+1, j+1} + (1+2r\lambda) u_{i+1, j} - \lambda r u_{i+1, j-1} = = r(1-\lambda) u_{i, j+1} + [1-2r(1-\lambda)] u_{i, j} + r(1-\lambda) u_{i, j-1},
\end{verbatim}

де вводиться позначення r = k/h\^{}2 . Зверніть увагу, що якщо i
фіксоване, а j змінюється від 2 до n-1, співвідношення (39.4) визначають
систему рівнянь n-2 з n-2 невідомим \texttt{u\_\{t+1,2\}},
\texttt{u\_\{t+1,3\}},

\pandocbounded{\includegraphics[keepaspectratio]{images/_page_296_Figure_5.jpeg}}

Фіг. 39.1. Сітка для схеми неявних різниць.

\pandocbounded{\includegraphics[keepaspectratio]{images/_page_296_Figure_7.jpeg}}

\begin{verbatim}
Фіг. 39.2. Шаблон неявної схеми u_{i+1, 4}, \ldots, u_{i+1, n-1} , які є розв'язком проблеми у внутрішнійMesh-вузли на часовому рівні t=(i+1)\Delta t .
\end{verbatim}

\begin{verbatim}
Фіг. 39.2 дає візуальне представлення структури кожного рівняння системи (39.4). Тепер перейдемо до розв'язку системи (39.4).# Алгоритм розв'язання задачі (39.2)КРОК 1. Виберіть якусь цінність \lambda ( 0 \le \lambda \le 1 ). Якщо \lambda = 0 , то рівняння (39.4) перетворюються на явні формули з лекції 38.КРОК 2. Припустимо, наприклад, h=\Delta x=0,2 і k=\Delta t=0,08 (у цьому випадку r=k/h^2=2 ). У цьому випадку сітка містить 6 вузлів уздовж осі (4 внутрішні вузли) (див. рисунок 39.1). Візьмемо параметр ваги \lambda=0,5 (отримана схема називається схемою Кренка-Ніколсона). Відповідно до обчислювальної схеми (її зазвичай називають шаблоном), показаної на рис. На рисунку 39.2, рухаючись зліва направо (j=2,3,4,5) уздовж перших двох шарів (i=1), отримуємо наступні чотири рівняння:
\end{verbatim}

\begin{verbatim}
\begin{array}{l} -u_{21}+3u_{22}-u_{23}=u_{11}-u_{12}+u_{13}=1,\\ -u_{22}+3u_{23}-u_{24}=u_{12}-u_{13}+u_{14}=1,\\ -u_{23}+3u_{24}-u_{25}=u_{13}-u_{14}+u_{15}=1,\\ -u_{24}+3u_{25}-u_{26}=u_{14}-u_{15}+u_{16}=1. \end{array}
\end{verbatim}

Перепишем их в матричной форме(39.5)

\begin{verbatim}
\begin{bmatrix} 3 & -1 & 0 & 0 \\ -1 & 3 & -1 & 0 \\ 0 & -1 & 3 & -1 \\ 0 & 0 & -1 & 3 \end{bmatrix}
\begin{bmatrix} u_{2\bar{z}} \\ u_{\bar{z}a} \\ u_{24} \\ u_{\bar{z}x} \end{bmatrix}
=
\begin{bmatrix} 1 \\ 1 \\ 1 \\ 1 \end{bmatrix}
\end{verbatim}

.

Матриця цієї системи називається тридіагональною. Щоб розв'язати систему
з трьома діагоналями

\begin{verbatim}
\begin{bmatrix} b_1 & c_1 & 0 & 0 & \dots & & & 0 \\ a_1 & b_2 & c_2 & 0 & \dots & & & 0 \\ 0 & a_{\bar{z}} & b_3 & c_3 & \dots & & & 0 \\ & & & & & & & & \\ \vdots & & & & & & & & \\ 0 & 0 & 0 & \dots & & & & a_{n-1} & b_n \end{bmatrix}
\begin{bmatrix} x_1 \\ x_{\bar{z}} \\ x_{\bar{s}} \\ \vdots \\ x_n \end{bmatrix}
=
\begin{bmatrix} d_t \\ d_{\bar{z}} \\ d_{\bar{s}} \\ \vdots \\ \vdots \\ d_n \end{bmatrix}
\end{verbatim}

,

преобразуем ее в эквивалентную систему вида

\begin{verbatim}
\begin{bmatrix} 1 & c_1^* & 0 & 0 & \dots & 0 \\ 0 & 1 & c_2^* & 0 & \dots & 0 \\ 0 & 0 & 1 & c_2^* & \dots & 0 \\ \vdots & \vdots & & & \vdots & \vdots \\ 0 & 0 & \dots & & 1 \end{bmatrix}
\begin{bmatrix} x_1 \\ x_2 \\ x_3 \\ \vdots \\ x_n \end{bmatrix}
=
\begin{bmatrix} d_1^* \\ d_2^* \\ d_3^* \\ \vdots \\ d_n^* \end{bmatrix}
\end{verbatim}

.

где

\begin{verbatim}
c_1^* = c_1/b_1
\end{verbatim}

\begin{verbatim}
,  c_{j+1}^* = \frac{c_{j+1}}{b_{j+1} - a_j c_j^*}  ,  j = 1, 2, \ldots, n-2  , и  d_1^* = d_1/b_1  ,  d_{j+1}^* = \frac{d_{j+1} - a_j d_j^*}{b_{j+1} - a_j c_j^*}  ,  j = 1,  2, \ldots, n-1  .У цьому перетворенні немає нічого дивного, оскільки друга система повністю еквівалентна першої. Матриця нової системи організована так, що цю систему дуже легко розв'язати. Розв'язуючи рівняння послідовно знизу вгору, отримуємо
\end{verbatim}

\begin{verbatim}
x_n = d_n^*, \quad x_j = d_j^* - c_j^* x_{j+1}, \quad j = n-1, \ n-2, \ \dots, \ 2, \ 1.
\end{verbatim}

\begin{verbatim}
Застосовуючи цей метод до системи чотирьох рівнянь (39,5), отримуємо розв'язок  u_{22}=0,60  ;  u_{23}=0,80  ;  u_{24}=0,80  ;  u_{23}=0,60  . Ми отримали наближені значення розв'язку у внутрішніх точках сітки на  t=\Delta t  . Тепер ми можемо зробити наступний крок у часі, але для цього нам доведеться розв'язати нову систему рівнянь.У неявній схемі кількість обчислень на кожному кроці більша, ніж на явному, але хорошу точність можна отримати навіть при значно більшому кроці.
\end{verbatim}

\section{ЗАВДАННЯ - 1. З диференціального рівняння (39.3) отримуємо
рівняння
(39.4).}\label{ux437ux430ux432ux434ux430ux43dux43dux44f---1.-ux437-ux434ux438ux444ux435ux440ux435ux43dux446ux456ux430ux43bux44cux43dux43eux433ux43e-ux440ux456ux432ux43dux44fux43dux43dux44f-39.3-ux43eux442ux440ux438ux43cux443ux454ux43cux43e-ux440ux456ux432ux43dux44fux43dux43dux44f-39.4.}

\begin{itemize}
\tightlist
\item
  \begin{enumerate}
  \def\labelenumi{\arabic{enumi}.}
  \setcounter{enumi}{1}
  \tightlist
  \item
    Побудуємо неявну схему скінченної дельти для задачі
  \end{enumerate}
\end{itemize}

\begin{verbatim}
\begin{array}{ll} (\mathrm{Y}\mathrm{H}\Pi) & u_t = u_{xx}, \quad 0 < x < 1, \quad 0 < t < \infty, \\ (\Gamma\mathrm{Y}) & \left\{ \begin{array}{ll} u_{\,0}(0, \ t) = 1, \\ u_{\,x}(1, \ t) + u_{\,}(1, \ t) = g_{\,}(t), \\ u_{\,x}(1, \ t) = 0, \quad 0 \leqslant \kappa \leqslant 1. \end{array}
\end{verbatim}

\begin{verbatim}
\right. & 0 < t < \infty, \end{array}
\end{verbatim}

\begin{enumerate}
\def\labelenumi{\arabic{enumi}.}
\setcounter{enumi}{2}
\tightlist
\item
  Как с помощью неявной схемы решить задачу(УЧП)
\end{enumerate}

\begin{verbatim}
u_t = u_{xx} + u, \quad 0 < x < 1,
\end{verbatim}

(ГУ) \$\$\$
\texttt{-\ 4.\ Как\ выглядит\ шаблон\ для\ уравнения\ (39.4),\ когда}
\lambda = 1 ? - 5. Постройте блок-схему решения уравнения
теплопроводности (39.2). Если есть возможность, то напишите программу
для ЭВМ Было бы очень полезно решить эту задачу численно при различных
значениях λ с начальным условием и (x, 0) = sin (πx). Сравните
полученные при различных λ численные результаты с аналитическим
решением, которое в данном случае имеет вид

\end{document}
